\subsection{PIQP for lifted interval constraints}\label{sec:method-piqp}
PIQP combines an infeasible-start interior-point method with the proximal
method of multipliers \cite{piqp}. The SDK sends its sparse interface
the convex lifted Hessian \(P_L\), affine equality block \(Ez=e\), and
remaining bounds \(\ell\le Cz\le u\). The lifted \(z\) comprises the
original decision and the costed side-specific or group-peak
auxiliaries. For the following derivation express the finite interval
sides as \(Gz\le h\), introducing a slack \(s\ge0\). A finite upper
side contributes \((C_i,u_i)\), a finite lower side
\((-C_i,-\ell_i)\). This conversion describes the mathematics; the SDK
passes the equality/interval representation rather than rebuilding
the solver's internal side coordinates.

Let \(y\) be the free equality multiplier, \(v\ge0\) the inequality
multiplier and \((\bar z,\bar y,\bar v)\) the proximal centers. In the
generic proximal system with positive parameters \(\rho,\delta\),
the stationary residuals and central target have the form
\begin{align}
 P_Lz+q_L+E^\top y+G^\top v+\rho(z-\bar z)&=0,\label{eq:piqp-rd}\\
 Ez-e-\delta(y-\bar y)&=0,\label{eq:piqp-re}\\
 Gz-h+s-\delta(v-\bar v)&=0,\qquad s_i v_i=\mu>0.\label{eq:piqp-ri}
\end{align}
Differentiating the last equality gives
\(v_i\Delta s_i+s_i\Delta v_i\) in the Newton system.
For \(\mu\) tending to zero and proximal centers agreeing with their
iterates, the shifted residuals reduce to the ordinary KKT conditions.
Positive proximal regularization changes the systems solved along the
path, which is useful for rank-deficient equality blocks. It does not
justify accepting a nonconvex input or treating a finite-iteration
termination as a proof. PIQP uses a predictor step, a centrality
correction, and a sparse \(LDL^\top\) solve of the regularized KKT
system. Sparsity and fill, the number of factorizations and iterations,
and finite-side count govern cost.

\begin{algorithm}[ht]
\caption{PIQP on the sparse lifted interval representation}\label{alg:piqp-route}
\begin{algorithmic}[1]
\State Assemble the exact lift and partition equality and interval rows.
\State Invoke the pinned sparse PIQP solver on a fresh instance.
\State Check the native termination and the dimensions of primal and dual output.
\If{status is \texttt{PIQP\_SOLVED}}
 \State Restore interval and equality multipliers in lifted row order.
 \State Recover the original decision and apply original-space and lifted checks.
\Else
 \State Record the native failure without an accepted primal or multiplier.
\EndIf
\end{algorithmic}
\end{algorithm}

The candidate route pins PIQP 0.6.4, absolute and relative tolerances
\(10^{-7}\), and at most \(200\,000\) iterations. Its tight reference
configuration uses \(10^{-10}\) and \(10^6\) iterations.
The interval dual is upper minus lower. Equality multipliers are
copied without imposing a sign. The lift and the physical-coordinate
evaluation are shared with the other routes, and numerical
acceptance still requires their independent residual and objective
checks after the native solved status.
